\chapter{Literature Review}
\label{chap:lit_review}

This chapter reviews the literature underpinning the four themes established in Chapter~1.
Section~\ref{sec:lit_perception} covers perception for industrial components, spanning pose
estimation, template matching, semantic segmentation, and size estimation from RGB-D information.
Section~\ref{sec:lit_dualarm} covers dual-arm manipulation and the coordination, motion planning and task-allocation strategies that shared workspaces demand. Section~\ref{sec:lit_collision} covers collision detection and distance computation, and Section~\ref{sec:lit_planning} the motion planning methods that consume those distances. Section~\ref{sec:lit_ik} covers inverse kinematics for manipulator control, and Section~\ref{sec:lit_rl} reinforcement learning for single-agent and multi-agent manipulation. Section~\ref{sec:lit_gaps} consolidates the gaps that motivate the work of this thesis.

%% ============================================================
\section{Perception for Industrial Component Handling}\label{sec:lit_perception}
%% ============================================================

Recent advancements in computer vision techniques have garnered significant attention from
researchers, leading to substantial contributions in the fields of object detection and object pose
estimation. This section traces that progress as it bears on the specific problem of small, low-textured industrial components such as nuts and bolts.

%% ------------------------------------------------------------
\subsection{Object Detection and 6D Pose Estimation}\label{subsec:lit_pose}
%% ------------------------------------------------------------

Early work established that stereo geometry alone can recover object pose in clutter.
\citet{sumi2002object} used stereo vision to determine the 6-DOF pose of an object placed in a
cluttered scene, even when undergoing occlusion, demonstrating good accuracy in estimating the pose
and improving recognition rates. Their work nonetheless faces challenges related to computational
complexity and is limited to certain lighting conditions. Learned approaches subsequently improved robustness. \citet{hu2018segmentation} introduced a segmentation-driven approach for 6D pose estimation. By predicting 2D keypoint locations for visible object parts and combining them using confidence scores, the method achieved robust pose estimation and demonstrated state-of-the-art performance on challenging datasets such as Occluded-LINEMOD and YCB-Video. Its reliance on accurate segmentation, and its potential sensitivity to occlusions and initial conditions, remain limitations.

Surface texture is a crucial feature for object recognition, but many industrial components lack
distinctive textural patterns, which is the central difficulty in this domain.
\citet{wohlhart2015learning} proposed learning descriptors using a CNN to address this challenge.
Their approach demonstrated potential for recovering pose information but carried high computational
complexity and a strong reliance on the dataset. Moreover, the method's robustness was compromised
by lighting variation, occlusions, and background clutter, all of which affect descriptor
accuracy. \citet{back2020segmenting} proposed a synthetic data generation pipeline together with an RGB-D Fusion Mask R-CNN to address the segmentation of texture-less industrial components. By randomizing textures in the synthetic data, the focus shifted to shape-based features, and a confidence map estimator enhanced the model's ability to exploit depth information for improved segmentation. While this approach demonstrated promising results, the reliance on synthetic data for training raises concerns about domain adaptation to real-world scenarios. Self-supervised methods for segmentation and data extraction have also been developed for downstream tasks such as robotic grasping, although developing such deep learning models is extensive and time-consuming.

Shape-based object segmentation can be advantageous in scenarios with poor figure-ground contrast or
illumination, as demonstrated by \citet{borenstein2006shape}. While such a framework performs well
on the shapes it has been trained on, its ability to generalize to completely unseen shapes may be
limited. \citet{wohlkinger2011shape} instead viewed the classification task as a matching problem,
introducing a method that identifies similar 3D models from comparable viewpoints by analyzing
specific surface features. Although the method incorporates inter-view similarity, it may still
struggle with extreme variations in viewpoint. Objects viewed from angles significantly different
from those in the training set might not be matched accurately. The approach may also have
limitations in handling severe occlusions, where significant parts of the object are not visible in
the depth image, which affects the accuracy of matching and classification.

\begin{table}[htbp]
\centering
\caption{Object detection and 6D pose estimation approaches for low-textured components.}
\label{tab:lit_pose}
\resizebox{\textwidth}{!}{%
\begin{tabular}{lll}
\toprule
\textbf{Approach} & \textbf{Principle} & \textbf{Reported limitation} \\
\midrule
\citet{sumi2002object}        & Segment-based stereo vision      & Computational cost, lighting-dependent \\
\citet{hu2018segmentation}    & Segmentation-driven 2D keypoints & Depends on accurate segmentation, occlusion-sensitive \\
\citet{wohlhart2015learning}  & CNN descriptor learning          & High cost, dataset reliance, clutter-sensitive \\
\citet{back2020segmenting}    & RGB-D fusion + synthetic data    & Domain adaptation to real scenes \\
\citet{borenstein2006shape}   & Shape-guided segmentation        & Limited generalization to unseen shapes \\
\citet{wohlkinger2011shape}   & Inter-view shape matching        & Extreme viewpoint change, severe occlusion \\
\bottomrule
\end{tabular}%
}
\end{table}

%% ------------------------------------------------------------
\subsection{Template Matching and Bin Picking}\label{subsec:lit_template}
%% ------------------------------------------------------------

Most bin-picking applications have utilized the concept of matching a template of the expected
object in the scene to estimate its 3D pose, as this is simpler and does not require any previously
trained models. This methodology can be preferable in industry because the CAD models used as
templates are typically already available for the components being
handled.

\citet{liu2012fast} addressed fast object localization and pose estimation in heavy clutter for
robotic bin picking, establishing that pose hypotheses can be generated and verified efficiently
enough for industrial cycle times. \citet{guo2020fast} developed a fast and robust bin-picking
system for densely piled industrial objects, targeting precisely the dense-pile regime in which
individual instances heavily occlude one another. \citet{nguyen2022templates} revisited templates
for 3D object pose estimation, demonstrating generalization to new objects and robustness to
occlusions, and thereby showing that the template paradigm remains competitive with learned
alternatives when generalization to unseen instances is required. The practical difficulty is that conventional template matching searches the entire sample space to find a suitable pattern, which is
computationally expensive. Restricting the search volume using a prior segmentation of the RGB image
is what makes the approach tractable in practice~\citep{sk2024bolthandling}.

%% ------------------------------------------------------------
\subsection{Semantic Segmentation Architectures}\label{subsec:lit_segmentation}
%% ------------------------------------------------------------

Object segmentation is a fundamental technique in computer vision, aiming to partition an image into
distinct regions based on the class or category to which each pixel belongs. There are two ways to
achieve segmentation, instance segmentation and semantic segmentation~\citep{yin2022bridging}.
Instance segmentation would classify all bolts within a cluttered image and assign a unique colour
and label to each individual bolt, whereas semantic segmentation would classify the bolts and assign
a single label and colour to all bolts in the image. Where the objective is to sort components by
size rather than to distinguish individual instances, a unique label per object is not required, and
using semantic segmentation simplifies the process and reduces computational
complexity. \citet{mo2022review} give a comprehensive review of classical 2D image segmentation models based on deep learning. Encoder-decoder CNN architectures have performed best for classical image segmentation. \citet{long2015fully} introduced fully convolutional networks, establishing that
classification architectures could be adapted to dense per-pixel prediction.
\citet{noh2015learning} proposed a deconvolution network that learns the upsampling path explicitly
rather than relying on fixed interpolation. \citet{ronneberger2015unet} introduced U-Net, whose skip
connections between encoder and decoder preserve spatial detail that pooling would otherwise
discard, and the DeepLab family extended this line with atrous convolution and spatial pyramid
pooling for multi-scale context. Alongside these, PPLiteSeg with its STDC~I and~II backbones provides lightweight CNN architectures designed for real-time applications in resource-constrained environments, trading a measure of accuracy for inference speed. OCRNet leverages an object-contextual representation approach, demonstrating the potential benefits of exploring techniques beyond classical CNN-based architectures by explicitly modeling the relationship between a pixel and the object region it belongs to. The comparative evaluation of these families on a custom
industrial dataset, rather than on general-purpose benchmarks, is what determines which is
appropriate for a given component-handling task.

\begin{table}[htbp]
\centering
\caption{Semantic segmentation architecture families relevant to industrial component handling.}
\label{tab:lit_seg}
\resizebox{\textwidth}{!}{%
\begin{tabular}{lll}
\toprule
\textbf{Architecture} & \textbf{Key idea} & \textbf{Trade-off} \\
\midrule
FCN~\citep{long2015fully}            & Dense per-pixel prediction from classifiers & Coarse boundaries \\
DeconvNet~\citep{noh2015learning}    & Learned upsampling path                     & Added parameters \\
U-Net~\citep{ronneberger2015unet}    & Encoder-decoder skip connections            & Memory for high resolution \\
DeepLab family~\citep{chen2018encoder}                        & Atrous convolution, pyramid pooling         & Heavier inference \\
PPLiteSeg (STDC~I/II)~\citep{peng2022pp}                 & Lightweight real-time backbone              & Accuracy traded for speed \\
OCRNet~\citep{yuan2020object}                               & Object-contextual representation            & Beyond-CNN complexity \\
\bottomrule
\end{tabular}%
}
\end{table}

%% ------------------------------------------------------------
\subsection{Size Estimation from RGB-D Information}\label{subsec:lit_size}
%% ------------------------------------------------------------

Size estimation from 2D segmented images, leveraging classical vision techniques, is the component
that converts a segmentation into an actionable physical measurement. \citet{omahony2020deep} compared traditional classical vision methods with deep learning methods, highlighting the advantages of integrating both approaches and explaining how the combination can
enhance computer vision perception, leading to more robust and accurate results. Classical vision
techniques such as markers, edge detection, corner detection, and contour detection may not be
reliable for all situations. Under varying light conditions, shapes, and colors these methods can
under perform at localization, consequently resulting in faulty localization and size
estimation.

Deep-learning-based object detection can estimate object size from bounding box dimensions and
aspect ratio with some success, as demonstrated using RGB-D cameras for size
estimation~\citep{wang2017tree}. Applying classical vision to a deep-learning-based
\emph{segmentation} instead provides more accurate object boundaries than direct bounding boxes,
which leads to accurate size estimation and eliminates the influence of background clutter on the
measurement. This distinction matters for small components, where a bounding box drawn around an elongated, obliquely oriented bolt encloses a substantial fraction of background, and any dimension inferred from that box inherits the error. Deep learning has also been applied more broadly to industrial inspection, where anomaly detection in visual data supports quality control alongside the component-handling tasks considered here~\citep{shukla2025anomaly}.

%% ============================================================
\section{Dual-Arm Manipulation and Coordination}\label{sec:lit_dualarm}
%% ============================================================

Dual-arm manipulators, mirroring bi-manual human configurations, have drawn attention for sorting
and object-transport applications~\citep{smith2012dual,makris2017dual}. Coordination between them
has been achieved by exchanging real-time kinematic data~\citep{ortenzi2018dual}, by sharing goal
configurations over a wireless link, or by relying on a global motion-capture system for centralized
spatial tracking~\citep{sk2025sharedvision}. Dual-arm systems can replace human workers in existing workspaces designed for bi-manual manipulation without requiring redesign~\citep{sk2025distance}. They have been proposed across a range of industrial scenarios, including complex welding tasks~\citep{gan2019multi}, agricultural harvesting~\citep{sepulveda2020robotic}, and peg-in-hole assembly with dexterous hands~\citep{lee2022peg}.

%% ------------------------------------------------------------
\subsection{Control-Level Coupling}\label{subsec:lit_coupling}
%% ------------------------------------------------------------

Early work coupled the arms at the control level, with both acting on one shared object.
\citet{uchiyama1988symmetric} proposed a symmetric hybrid position/force control scheme for the
coordination of two robots, in which workspace force, velocity and position vectors are symmetric
functions of both arms' joint-space quantities, so that neither arm dominates. The leader-follower
alternative instead commands one arm's trajectory directly and derives the second arm's motion to
accommodate it, and these and related hybrid and adaptive variants are surveyed by
\citet{caccavale2025cooperative} and \citet{smith2012dual}. Throughout this line of work,
coordination concerns how two arms move together on the same object, not which arm should
act.

\citet{abbas2023systematic} provide a systematic review of cooperative dual-arm manipulators
covering modelling, planning, control, and vision strategies, and \citet{makris2021cooperative}
categorizes cooperative manipulation into non-coordinated manipulations, where the manipulators move
independently executing different tasks, and coordinated manipulations, involving combined movements
with spatial and temporal synchronization in a shared workspace. Coordinated manipulation divides
further into symmetric, with mirror movements and equal force distribution, and asymmetric, with
complementary movements and different force distributions. Although more challenging, asymmetric
manipulation enables a wider range of complex tasks that require coordinated
actions~\citep{surya2025iapf}. Learned approaches to the same coupling problem have also appeared. \citet{alles2022learning} address learning to centralize dual-arm assembly, and \citet{zhang2025imagebased} develop image-based visual servoing for enhanced cooperation of dual-arm manipulation. More recently, \citet{chen2026vlmsfd} propose a vision-language-model-assisted flow diffusion framework for dual-arm cooperative manipulation. Eye-in-hand vision is well established for single-arm pick-and-place~\citep{nguyen2025vision}, and dual-arm systems have been used for perception-intensive tasks such as cooperative 3D model acquisition~\citep{kobayashi2022obtaining}, although there the sensing comes from a stationary external depth sensor rather than each arm's own wrist camera. In neither case is a second arm's camera used to decide which arm should manipulate a shared target.

%% ------------------------------------------------------------
\subsection{Task Allocation in Multi-Robot and Dual-Arm Systems}\label{subsec:lit_allocation}
%% ------------------------------------------------------------

Efficient task allocation in robot systems enhances overall performance and resolves workspace
reachability constraints. \citet{khamis2015multi} review the state of the art in multi-robot task
allocation, and \citet{korsah2013comprehensive} provide a comprehensive taxonomy for the problem.
\citet{alatartsev2015robotic} survey the closely related robotic task sequencing problem, and
\citet{keshvarparast2024collaborative} review collaborative robots in manufacturing and assembly
systems together with a future research agenda. \citet{pedersen2016robot} address the translation of
robot skills for manufacturing from concept to industrial deployment, and \citet{patra2024online}
develop online capability-based task allocation of cooperative manipulators.

Market-based task allocation methods have emerged as a prominent distributed approach, in which
robots bid for tasks based on their reachability and current
states~\citep{quinton2023market,galati2023auction}. \citet{fazal2022task} proposed a dynamic
threshold-based system using adaptive incentive value computation with sequential single-item
auctions and MiniSum/MiniMax optimization. However, centralized incentive computation and continuous
system-wide resource tracking create computational bottlenecks unsuitable for real-time
applications. \citet{braquet2021greedy} proposed a greedy coalition auction algorithm for spatially
distributed multi-agent systems, enabling dynamic task reassignment as agents move toward targets
through decentralized bid negotiations, and \citet{schneider2015auction} examine auction-based
allocation for multi-robot teams in dynamic environments. Round-robin task allocation provides
computational simplicity and straightforward assignment, but fails to account for task complexity
and spatial conflict handling, leading to suboptimal utilization in spatially constrained
workspaces~\citep{alirezazadeh2024survey}.

Recent learning-based approaches employ deep reinforcement learning for adaptive task allocation and
have shown promise in multi-robot coordination~\citep{cui2024task, hao2024coordinated}, with DRL approaches achieving better performance than heuristic baselines through end-to-end policy
optimization~\citep{ma2024end,lu2024reinforcement}. These methods primarily address spatially
distributed scenarios in which robots operate in large-scale workspaces such as warehouses and
agricultural fields, learning task-scheduling decisions from abstract state representations such as
resource availability, geometric proximity, and task priorities. Centralized formulations schedule
tasks based on robots' current positions, predetermined target distances, and battery states, while
decentralized DQN-based approaches learn to make run, transfer, or postpone decisions from local
resource usage and task priority, notably excluding physical task locations or multi-robot spatial
relationships~\citep{gautier2020comparison}. Moreover, these learning-based methods require
extensive offline training and intensive iterations before taking a decision compared with reactive
methods.

In dual-arm systems, task allocation presents unique challenges due to shared workspace constraints,
kinematic coupling, and the need for coordinated bilateral manipulation in industrial environments.
\citet{behrens2019constraint} cast simultaneous task allocation and motion scheduling for industrial
dual-arm cells as a constraint program, jointly optimizing task order and arm assignment under
precedence and timing constraints. \citet{hamasaki2023autonomous} employed a three-layered
A*/anytime-A* architecture with constraint-based conflict resolution and stepwise rescheduling for
online disruptions. While achieving good performance in structured construction tasks, the
multi-layer planning architecture and constraint resolution approach may not be optimal for
high-frequency object detection scenarios requiring immediate proximity-based allocation decisions.
\citet{rodriguez2016combining} proposed hybrid motion and task planning using AND/OR manipulation
task graphs with motion planner feedback. While providing systematic solutions through backtracking
search, the approach requires complete workspace knowledge and pre-computation of grasp
configurations, making it unsuitable for dynamic real-time environments.
\citet{kimmel2016scheduling} addressed dual-arm coordination using coordination diagrams with batch
and incremental methods for task assignment. Although computationally efficient, both approaches
assume known object placements and require collision-free trajectory pre-computation, limiting
applicability to vision-based dynamic scenarios.
\citet{qin2024onlinemulti} address online task allocation and scheduling in multi-manipulator
systems under collision constraints and unknown tasks, and \citet{zacharias2007capturing} address
the assignment question at the design stage by building reachability and capability maps that
support a static, offline division of the workspace into per-arm zones.

\begin{table}[htbp]
\centering
\caption{Task allocation strategies for multi-robot and dual-arm systems.}
\label{tab:lit_allocation}
\resizebox{\textwidth}{!}{%
\begin{tabular}{lll}
\toprule
\textbf{Strategy} & \textbf{Mechanism} & \textbf{Reported limitation} \\
\midrule
Market-based~\citep{quinton2023market}   & Bidding on reachability and state & Centralized bottlenecks \\
Auction~\citep{braquet2021greedy}        & Decentralized bid negotiation      & Convergence under churn \\
Threshold~\citep{fazal2022task}          & Adaptive incentive value           & System-wide resource tracking \\
Round-robin~\citep{korsah2013comprehensive} & Fixed cyclic assignment         & Ignores spatial conflict \\
Constraint program~\citep{behrens2019constraint} & Joint order and assignment & Offline, needs full model \\
Layered A*~\citep{hamasaki2023autonomous} & Search with rescheduling          & Not for high-frequency detection \\
Coordination diagrams~\citep{kimmel2016scheduling} & Incremental search       & Assumes known object placement \\
Deep RL~\citep{ma2024end}                & End-to-end policy                  & Extensive offline training \\
\bottomrule
\end{tabular}%
}
\end{table}

%% ============================================================
\section{Collision Detection and Distance Computation}\label{sec:lit_collision}
%% ============================================================

Industrial automation adopted multi-arm manipulator systems driven by the need for fast production,
and researchers focused on the development of collision detection methods as multi-arm robotic
systems emerged as a critical research area. Collision detection studies in robotics have ranged
from geometric methods to learning-based methods.

%% ------------------------------------------------------------
\subsection{Geometric Methods}\label{subsec:lit_geometric}
%% ------------------------------------------------------------

\citet{jimenez2001collision} presented a comprehensive review of collision detection methods in 3D
space, and \citet{lin1993efficient} established foundational techniques for efficient collision
detection in animation and robotics. \citet{gilbert1988fast} introduced a fast procedure for
computing the distance between complex objects in three-dimensional space, a formulation whose
acceleration continues to attract attention~\citep{montaut2024gjk}. \citet{redon2005fast} addressed
fast continuous collision detection for articulated models, extending the problem from discrete
configurations to continuous motion between them.

\citet{lumelsky1985fast} introduced a parameter-based algorithm for minimum distance calculation
between line segments in $n$-dimensional space. Using parameters that range over the unit interval
for points lying inside the segments, the method establishes a framework handling all geometric
configurations. It first calculates the minimum distance between infinite lines by minimizing a
quadratic function, then verifies whether these points lie within the actual segments. By using
parameter values to identify the relevant endpoints, computational complexity reduces to $5n+12$
multiplications and $8n+1$ additions. However, it provides only discrete collision detection without
continuous distance metrics, limiting its application in modern control approaches requiring
gradient information for smooth motion planning.

Subsequent geometric work refined the primitives. Collision detection approaches using minimum
distance functions defined by Euclidean norms determine collision risk between robot arms by
modeling robot links as cylinders with hemispheres and the end-effector as spheres, then
calculating distances between these primitives in three scenarios, line segment to line segment,
line segment to point, and point to point~\citep{chang1990collision}. Hierarchical geometric collision detection algorithms for dual-arms have combined bounding boxes, joint-joint distance thresholds, perpendicular distance calculations, and cross-product methods~\citep{lee2001real}. Despite a 30\% computational improvement, such schemes used simplified primitives, lacked mathematical formalism for complex configurations, and provided only discrete binary detection without continuous metrics for smooth planning. These approaches showed reduced accuracy with near-parallel links and increased complexity in 3-D, without addressing degenerate cases or providing control gradient information.

Dual-stage collision detection algorithms~\cite{ketchel2008self}, for spatial closed chains first use dual vector algebra and Pl\"ucker~\citep{miller2005plucker} coordinates to quickly check for possible collisions between infinite cylinders, and, if potential collisions are detected, perform precise second-stage testing between finite cylinders using detailed geometric calculations. While mathematically elegant, this approach may not provide the continuous distance metrics needed by modern control systems for obstacle avoidance, and it formulates collision detection as a constrained nonlinear minimization problem, which introduces complexity for real-time applications by requiring conjugate gradient methods with barrier functions to solve. \citet{larsen2000fast} presented an RSS-based proximity algorithm that combines rectangles with uniform sphere expansion to create efficient bounding volumes, building hierarchical trees using covariance analysis and geometric slab properties for fast distance calculations. Performance improvements come from priority-directed search and triangle caching that leverages frame-to-frame coherence. While outperforming sphere-based hierarchies, the approach struggles with non-convex geometry, requiring decomposition into multiple parts, and can experience numerical instabilities in near-intersection cases, sometimes producing overly conservative distance estimates that reduce query efficiency.

More recent formulations pursue continuous metrics directly. \citet{xu2020pseudo} proposed a
pseudo-distance algorithm for collision detection of manipulators using a
convex-plane-polygons-based representation, and \citet{xia2024collision} formulate collision
detection between convex objects using pseudo-distance and unconstrained optimization. The necessity
of collision detection methods is not limited to manipulators but extends to autonomous vehicle
motion planning and multi-agent path planning, and the human-robot case has been addressed directly
through multi-level optimization for fast distance calculation between a human and a manipulator,
and through early prediction of human motion for collaborative manipulation
planning~\citep{mainprice2013human}.

\begin{table}[htbp]
\centering
\caption{Collision detection and distance computation methods, classified by the form of output they
provide to a controller.}
\label{tab:lit_collision}
\resizebox{\textwidth}{!}{%
\begin{tabular}{llll}
\toprule
\textbf{Method} & \textbf{Representation} & \textbf{Output} & \textbf{Limitation} \\
\midrule
\citet{lumelsky1985fast}   & Line segments        & Discrete distance   & No gradient for smooth planning \\
\citet{gilbert1988fast}    & Convex polytopes     & Distance            & Convexity assumption \\
\citet{larsen2000fast}     & Swept-sphere volumes & Bounded distance    & Non-convex geometry, instability \\
\citet{redon2005fast}      & Articulated models   & Continuous contact  & Cost for many links \\
\citet{xu2020pseudo}       & Convex plane polygons& Pseudo-distance     & Representation build cost \\
\citet{montaut2024gjk}     & Accelerated GJK      & Distance            & Convexity assumption \\
\citet{das2020learning}    & Data-driven          & Binary / threshold  & Black box, large data requirement \\
\bottomrule
\end{tabular}%
}
\end{table}

%% ------------------------------------------------------------
\subsection{Learning-Based Methods}\label{subsec:lit_learned_collision}
%% ------------------------------------------------------------

Learning-based methods have attracted attention from researchers due to their potential to
accelerate collision detection applications. Machine-learning models have been introduced that serve
as better alternatives to conventional geometric methods, achieving considerable speedup in motion
planning applications. In the same way, the minimum distance between manipulator links has been used
as a criterion for collision detection, with a neural network trained on distance data for collision
prediction against a fixed threshold~\citep{das2020learning}. While learning methods can learn complex non-linear collision patterns from labeled data, the learning mechanisms are black boxes and need large amounts of data for learning, which limits their applicability where interpretability or guaranteed behavior is required. This trade-off, between the interpretability and determinism of geometric methods and the speed of learned ones, is what motivates a closed-form analytical treatment that retains continuous distance metrics while remaining cheap enough for real-time evaluation.

%% ============================================================
\section{Motion Planning and Artificial Potential Fields}\label{sec:lit_planning}
%% ============================================================

The planning of dual-arm movements addresses coordination through two primary approaches. Path
planning generates geometric trajectories avoiding obstacles and self-collisions, prioritizing
spatial relationships without timing information. Motion planning extends this by incorporating
velocity profiles, acceleration constraints, and dynamic considerations, creating time-parametrized
trajectories~\citep{abbas2023systematic}. The selection between these approaches depends on
application requirements, with industrial manipulation tasks often benefiting from motion planning's
ability to manage the temporal aspects of coordination between manipulators. \citet{tang2012review} provide a comprehensive review of robot motion planning approaches, categorizing and analysing the research evolution from 1980 to 2010. They systematically compared conventional approaches such as Bug algorithms, roadmaps, cell decomposition, potential fields, and mathematical programming with heuristic methods including neural networks, genetic algorithms, and fuzzy logic, demonstrating the field's evolution from conventional methods in the 1980s to predominantly heuristic approaches by the 2000s. Their analysis includes quantitative comparisons of implementation frequency for each approach and identifies future research directions for integrating perception, sensing, and motion planning.

\citet{khatib1986real} pioneered the Artificial Potential Field method, transforming discrete path
planning into a continuous optimization problem in which the robot is treated as a particle moving
under the influence of attractive forces toward the goal and repulsive forces away from obstacles.
The formulation's appeal is that it is reactive and cheap to evaluate, producing a control signal
directly rather than a path that must subsequently be tracked. \citet{volpe1987artificial} refined
the original concept by introducing elliptical iso-potential contours that vary with distance,
improving convergence properties, and potential-field formulations have since been extended to
real-time multi-agent navigation in industrial environments~\citep{prajapati2026mppf}.

For dual-arm systems, however, APF methods face unique challenges related to workspace sharing and
coordinated manipulation in manufacturing contexts. Alternating master-slave paradigms with charged
surface workspace modeling have been employed for collision avoidance~\cite{chuang2006alternate}. However, the sequential trajectory planning approach lacks true coordination and may cause deadlocks in constrained workspaces. Exponential-field variants developed for aerial navigation with vertical motion prioritization suffer from oscillations in complex obstacle environments\citep{jayaweera2020dapf}. Hybrid combinations of APF with sampling-based planners have been proposed for mobile dual-arm systems, and force sensing has been integrated with rotating repulsive fields for enhanced navigation capabilities, but neither resolves the asymmetric dual-arm case~\citep{lin2025robot}. Hierarchical leader-follower deep reinforcement learning frameworks using Soft Actor-Critic have been proposed for dual-arm closed-chain manipulation in cooperative transportation\citep{cui2024task, hao2024coordinated}. While such approaches demonstrate task adaptability to random target poses, they require extensive offline training for convergence, and the fixed leader-follower structure limits applicability, performing well in symmetric manipulation scenarios but struggling with asymmetric dual-arm coordination. Once task allocation is determined, the fundamental challenge shifts to coordinated motion execution while ensuring collision-free trajectories in shared workspace environments, which is where the continuous distance metrics of Section~\ref{sec:lit_collision} become the input to the planner rather than merely a safety monitor.

%% ============================================================
\section{Inverse Kinematics for Manipulator Control}\label{sec:lit_ik}
%% ============================================================

Robotic manipulators are indispensable to modern industrial automation, serving as the backbone for
tasks ranging from precision assembly to complex material handling. Traditionally, spherical wrist
designs have dominated the industry due to their kinematic decoupling, which facilitates independent
closed-form analytical solutions for position and
orientation~\citep{rodriguez2024singularity}. While these designs have long been the industrial
standard, they impose significant structural limitations~\citep{shah2019comparison}. The geometric
constraint requiring the final three axes to intersect at a common point often compromises
mechanical stiffness and restricts the integration of high-torque motors within compact wrist
assemblies~\citep{ostermeier2025automatic}. Consequently, non-spherical wrist manipulators have
gained prominence, achieving superior structural integrity and higher payload-to-weight ratios by
introducing offsets between the last three axes~\citep{carbonari2023inverse}. This violation of
kinematic decoupling renders analytical inverse kinematics intractable, necessitating reliance on
complex numerical or data-driven solvers.

%% ------------------------------------------------------------
\subsection{Analytical and Numerical Solvers}\label{subsec:lit_analytical}
%% ------------------------------------------------------------

To address the general inverse kinematics problem, several classical methods including numerical,
analytical, combined, and heuristic frameworks have been used in the literature. Iterative solvers
such as TRAC-IK~\citep{beeson2015trac} are widely used for their system-agnostic ability to handle
joint limits and multiple constraints. However, even with a previous joint-state seed and a distance
parameter for solution selection, TRAC-IK remains fundamentally trajectory-unaware. Each waypoint is
resolved as an independent optimization with no inter-waypoint coupling, leading to
configuration-branch switching and joint-space discontinuities during continuous trajectory
tracking. Heuristic methods such as FABRIK~\citep{aristidou2011fabrik} are computationally efficient but limited to position-reaching applications and lack the necessary temporal awareness. Analytical closed-form solutions reduce the complex kinematics of the robot into a single 16th-degree polynomial to calculate all possible joint configurations for a given pose~\citep{zohour2021kinova}. Without strict nearest-joint constraints, continuous sub-millisecond
tracking causes the solver to jump between completely different joint configurations. Automatic geometric decomposition has been proposed to recover analytical solutions where the structure permits~\citep{ostermeier2025automatic}, and metaheuristic algorithms have been compared for solving manipulator inverse kinematics, but neither removes the fundamental coupling introduced by wrist offsets.

%% ------------------------------------------------------------
\subsection{Velocity-Level Controllers}\label{subsec:lit_velocity}
%% ------------------------------------------------------------

Velocity-level controllers address the trajectory-awareness problem by operating directly in the
differential domain. Resolved-rate motion control~\citep{whitney1969resolved} computes joint
velocities using the Jacobian pseudo-inverse. Damped least squares
formulations~\citep{wampler1986manipulator,chiaverini1994review} introduce a damping factor to
prevent velocity singularities, trading a small tracking error near singular configurations for
bounded joint velocities, and \citet{nakamura1986inverse} established the singularity-robustness
properties that make this practical for hardware. Quadratic velocity-level programming
formulations~\citep{zhang2013variable} further impose joint limit constraints explicitly, offering
safer hardware deployment, and null-space saturation methods handle joint constraints in redundant
manipulators~\citep{flacco2012motion}. Model predictive contouring control has been applied to
manipulators for reactive trajectory following~\citep{yoon2025reactive}. All of these methods,
however, require an explicit, accurate Jacobian at every timestep, which is the cost that
learning-based differential formulations seek to remove. Visual servo
control provides the complementary sensing side of this loop, computing Cartesian pose error
directly from image measurements~\citep{hutchinson1996tutorial}.

%% ------------------------------------------------------------
\subsection{Learning-Based and Hybrid Approaches}\label{subsec:lit_learned_ik}
%% ------------------------------------------------------------

To improve upon the accuracy limitations of classical velocity controllers, hybrid approaches have
emerged combining numerical methods with learned components~\citep{wang2023two}. Learning-based methods have been explored as a residual corrector atop numerical solvers, compensating for solver inaccuracies~\citep{jia2019finite, chen2019neural}. An actor-critic RL approach has been introduced as a correction signal in joint velocity space to an augmented PD controller for trajectory tracking on a 6-DOF industrial manipulator~\citep{pane2016actor}, and~\citet{aghaei2022real} a
gradient-free, model-free Bayesian Markov Chain Monte Carlo RL policy search has been used for
circular trajectory control on a 2-DOF manipulator in simulation to optimize a PD controller. A Soft
Actor-Critic based outer-loop stochastic reference correction policy has been learned directly on a
7-DOF manipulator, augmenting an existing vendor-provided hardware controller with bounded
corrective actions as an inheritance rather than learning from the environment~\citep{pavlichenko2022real}. Distributed PPO with
LSTM and generalized advantage estimation has been applied to trajectory tracking on lower-DOF
platforms using random reference state initialization and an exponential position-error
reward~\citep{zhang2019trajectory}.

Compound controllers have also been proposed, combining a sliding mode controller with SAC-based RL
augmented by Lyapunov stability constraints for manipulator trajectory tracking~\citep{liu2022reinforcement, mao2025real}. While this requires
a complete dynamic model of the system, it demonstrates that model-based initialization
significantly improves RL sample efficiency. A hybrid framework combining Twin Delayed Deep
Deterministic Policy Gradient with an RRT-generated reference trajectory penalizes deviation from
the reference directly, constraining the agent to refine a pre-computed joint-space path~\citep{ota2019trajectory}.~\citet{elumalai2025proximal} proposed a collaborative PPO variant has been formulated for tracking and vibration suppression of a flexible-joint manipulator, reporting that PPO-based policies converge faster and more stably than actor-critic and policy-gradient baselines under hardware-in-the-loop validation, and~\citet{li2021general} proposed traditional path planner-DRL frameworks have combined RRT* with actor-critic algorithms to find energy-optimized IK solutions for redundant manipulators, treating IK as a pose-reaching problem solved point by point.

Several DDPG-based methods have been explored to enhance trajectory tracking, though these often
depend on simplified assumptions or known dynamics.~\citet{majumder2024enhancing} proposed sequential DDPG approaches for $n$-link manipulators utilize dual actor-critic networks trained progressively from 2D to 3D space to accelerate convergence. Distributed DDPG frameworks with prioritized experience replay employ multiple synchronous workers that explore diverse policies to update a shared global network~\citep{zhang2019trajectory}. Hybrid variants learn joint torque commands by embedding an explicitly derived Euler-Lagrange dynamic model within the simulation~\citep{hazem2025reinforcement}. While these methods demonstrate efficacy in simulation, they frequently operate in joint torque space requiring precise manipulator dynamics, and lack real-time hardware validation, leaving the sim-to-real gap unaddressed. Moreover, these frameworks often overlook orientation control and motion smoothness, which are critical for scaling to 6-DOF manipulators.

\begin{table}[htbp]
\centering
\caption{Inverse kinematics solver families and their suitability for continuous trajectory tracking
on a non-spherical wrist.}
\label{tab:lit_ik}
\resizebox{\textwidth}{!}{%
\begin{tabular}{llll}
\toprule
\textbf{Family} & \textbf{Representative} & \textbf{Trajectory aware} & \textbf{Requirement} \\
\midrule
Iterative numerical & TRAC-IK~\citep{beeson2015trac}          & \xmark & Per-waypoint optimization \\
Heuristic           & FABRIK~\citep{aristidou2011fabrik}      & \xmark & Position only \\
Analytical          & Polynomial~\citep{zohour2021kinova}     & \xmark & Nearest-joint constraint \\
Resolved-rate       & RRMC~\citep{whitney1969resolved}        & \cmark & Explicit Jacobian \\
Damped least squares& DLS~\citep{chiaverini1994review}        & \cmark & Explicit Jacobian, damping \\
Quadratic program   & QP~\citep{zhang2013variable}            & \cmark & Explicit Jacobian, solver \\
Learned differential& DIKPPO~\citep{sk2025diktracking}           & \cmark & Training, no Jacobian \\
\bottomrule
\end{tabular}%
}
\end{table}

%% ============================================================
\section{Reinforcement Learning for Manipulation}\label{sec:lit_rl}
%% ============================================================

Intelligent robotic manipulation is a key element of autonomous operations in industrial automation
systems. Recent advancements in robot learning have significantly accelerated the pace of this
research, especially through the integration of Imitation Learning and Reinforcement Learning for
the autonomous generation of complex and intelligent robotic manipulation
skills~\citep{zare2024survey}.

%% ------------------------------------------------------------
\subsection{Single-Agent Reinforcement Learning}\label{subsec:lit_single_rl}
%% ------------------------------------------------------------

Conventional robotic control methods typically involve modular pipelines consisting of sequential
planning and control stages. These methods depend on accurate kinematic or dynamic models, which are
often complex and difficult to obtain in practice, and predefined control sequences struggle to
generalize across task variations without explicit re-engineering of each module~\citep{he2023robotic}. Data-driven control techniques such as Learning from Demonstration, and data collection via teach pendants or tele-operation, do not require explicit knowledge of robot dynamics or environmental models to execute tasks. Although these approaches eliminate the dependency on accurate dynamic models, they require a large volume of high-quality data, so that performance remains constrained by the availability and diversity of the collected datasets, often making them difficult to scale in complex industrial environments~\citep{zhang2015towards}.

To overcome these limitations of data scarcity and human dependence on data collection,
reinforcement learning has been introduced as a powerful
alternative~\citep{kober2013reinforcement}. Unlike data-dependent IL and LfD, RL enables a robot to
figure out optimal control policies through interaction and exploration with the surroundings.
Although RL methods require extensive training time to converge, their capacity to learn optimal
control robustly is often more reliable than static data-driven approaches, and by optimizing a
well-defined reward function, RL-based agents can potentially surpass the performance of human
experts and develop resilient recovery behaviors that are not represented in static
datasets~\citep{zare2024survey, brune2025deep, gu2017deep}.

Recent studies demonstrate that RL is highly effective for robot control across various domains,
including trajectory optimization~\citep{majumder2024enhancing, mao2025real}, intelligent pick-and-place operations~\citep{deng2019deep, mehta2024waypoint}, obstacle avoidance~\citep{zhu2022robot, al2024reinforcement}, and human-robot interaction tasks~\citep{qureshi2018intrinsically, zhang2022reinforcement}. RL enables robots to learn from experience and provides a flexible framework for handling non-linear dynamics and unstructured environments that often challenge traditional control methodologies. However, extending RL to multi-agent configurations poses a high-dimensional action space that leads to exponential growth in exploration complexity~\citep{low2025cooperative, nguyen2020deep}. Furthermore, the requirement for precise spatio-temporal synchronization between agents such as a dual-arm system makes the reward assignment problem significantly more challenging than in simple single-agent applications.

%% ------------------------------------------------------------
\subsection{Multi-Agent Reinforcement Learning and CTDE}\label{subsec:lit_marl}
%% ------------------------------------------------------------

Although single-arm manipulation based on RL has demonstrated promising results, complex tasks such
as alignment and assembly of components require collaborative manipulation rather than isolated
single-arm execution. Extending RL from single-agent to multi-agent systems poses significant
challenges. In a single-agent environment, an agent is concerned only with the outcomes of its own
actions within a stationary distribution. However, in multi-agent domains, an agent must account for
both the consequences of its own actions and the evolving behaviors of other agents. This
interaction introduces non-stationarity into the learning process, as the optimal policy of one
agent becomes dependent on the fluctuating strategies of its
counterparts~\citep{low2025cooperative,nguyen2020deep}.

Multi-Agent Reinforcement Learning was introduced to address these challenges by explicitly
modelling the interactions between multiple decision-making
entities~\citep{albrecht2024multi,low2025cooperative}. MARL methods are broadly classified along
three primary axes. First, by cooperation structure. Agents may be fully cooperative, sharing a
joint reward to achieve a common goal, fully competitive, operating in zero-sum settings, or mixed,
pursuing partially conflicting objectives. Second, by the training and execution paradigm.
Independent Learning treats other agents as part of a non-stationary
environment~\citep{dewitt2020independent}, Centralized Training with Decentralized Execution enables
agents to learn a shared critic during training while maintaining independent policies during
execution, and Centralized Learning subordinates all agents to a single controlling entity. Third,
by communication protocol. Communicative agents explicitly share state or intent information, while
non-communicative agents must implicitly infer the behaviors of others through observed
actions~\citep{liu2024efficient, chai2025survey}.

For dual-arm manipulation, the problem is most effectively formulated as a fully cooperative
decentralized partially observable Markov Decision Process under the CTDE
paradigm~\citep{yu2022surprising}. While recent heterogeneous-agent CTDE methods such as HATRPO and
HAPPO~\citep{} offer monotonic improvement guarantees for asymmetric tasks via sequential policy
updates~\citep{kuba2021trust}, agents in a symmetric dual-arm system are structurally symmetric,
each arm sharing identical architecture and observation dimensionality. In this homogeneous
cooperative setting, MAPPO's simultaneous update with a shared global-state critic is well suited,
directly capturing inter-arm spatial relationships during training while preserving decentralized
execution at deployment~\citep{chai2025survey}.

Despite the theoretical advantages of CTDE frameworks, their direct application to dual-arm
collaborative manipulation presents significant practical problems. Specifically, when agents are
required to simultaneously acquire low-level actuator skills and high-level inter-agent coordination
from scratch, sample efficiency degrades exponentially~\citep{yu2022surprising, zhu2020ingredients}. Current dual-arm MARL approaches often rely on rigid hand-designed action priors that do not adapt to different task phases. Furthermore, the existing literature lacks a principled mechanism for smoothly transitioning between stable single-arm priors and learned multi-agent coordination policies. This leaves the entire computational burden on the MARL agents, ultimately limiting the scalability and sim-to-real transferability of the learned controllers.

%% ============================================================
\section{Research Gaps}\label{sec:lit_gaps}
%% ============================================================

The review above consolidates into seven gaps, which correspond to those introduced in Chapter~1 and
are restated here in the terms established by this chapter.

\begin{enumerate}
    \item \textbf{Perception for small, low-feature components.} Pose estimation methods depend on
    accurate segmentation and degrade under occlusion and lighting variation. Descriptor learning
    carries high computational cost and dataset dependence, and synthetic-data approaches raise
    domain adaptation concerns. Existing sorting research remains limited in handling densely packed
    objects with limited visual features, and many methods require multiple sensors that increase
    system complexity and cost~\citep{sk2024bolthandling,sk2024boltsorting}.

    \item \textbf{Runtime arm selection from onboard perception.} Coordination has been addressed by
    control-level coupling, by centrally solved scheduling from an external view, or by offline
    static workspace zoning. In no case is a second arm's own eye-in-hand camera used to decide
    which arm should manipulate a shared target at runtime~\citep{sk2025sharedvision}.

    \item \textbf{Sensor-free, real-time inter-arm distance.} Conventional collision detection either
    depends on external sensors or on computationally expensive algorithms that compromise real-time
    applicability. Classical geometric formulations provide only discrete detection without the
    continuous, differentiable distance metrics that modern controllers require for smooth motion
    planning, while learned detectors are black boxes requiring large
    datasets~\citep{sk2025distance,surya2025iapf}.

    \item \textbf{Stable asymmetric dual-arm motion planning.} Asymmetric coordinated manipulation
    enables a wider range of complex tasks but is more challenging, and APF methods face unique
    difficulties related to workspace sharing, deadlock in constrained workspaces, and oscillation
    in complex obstacle environments~\citep{surya2025iapf,sk2025competitive}.

    \item \textbf{Adaptive task allocation.} Current industrial practice assigns fixed object types
    to each arm and fails when objects appear outside designated zones. Market-based methods incur
    centralized bottlenecks, round-robin ignores spatial conflict, and learning-based allocation
    requires extensive offline training while excluding physical task
    locations~\citep{sk2025competitive}.

    \item \textbf{Trajectory tracking on non-spherical wrists.} The violation of kinematic decoupling
    renders analytical IK intractable, iterative solvers remain trajectory-unaware, and classical
    velocity-level controllers require an explicit, accurate Jacobian at every timestep. No prior
    work demonstrates that a policy trained purely on static pose reaching, without trajectory
    supervision, generalizes to continuous trajectory tracking at
    deployment~\citep{sk2025diktracking}.

    \item \textbf{Sample-efficient dual-arm coordination learning.} When agents must simultaneously
    acquire low-level actuator skills and high-level coordination from scratch, sample efficiency
    degrades exponentially, and the literature lacks a principled mechanism for smoothly
    transitioning between stable single-arm priors and learned multi-agent coordination
    policies~\citep{sk2025alpha}.
\end{enumerate}

The chapters that follow address these gaps in the order established in Chapter~1.
